% Original proof clinic; mathematical background: CMU 15-317 natural deduction.
\documentclass[12pt,aspectratio=43]{beamer}
\usetheme{CambridgeUS}
\usecolortheme{dove}
\usefonttheme{serif}
\usepackage{amsmath,amssymb,bussproofs}
\setbeamertemplate{theorems}[numbered]
\setbeamertemplate{navigation symbols}{}
\title{Can you build the proof?}
\subtitle{A clinic in constructive natural deduction}
\author{Logic teaching fixtures}
\date{Seminar 2}
\newcommand{\seq}[2]{#1\vdash #2}
\begin{document}
\section{The challenge}
\begin{frame}{Start with a challenge, not a rule list}
\centering
\Large Given $A\land B$, construct $B\land A$.
\bigskip
\normalsize
\begin{enumerate}
\item Work alone for one minute.
\item Name every assumption you use.
\begin{itemize}\item What is given?\item What is the target?\end{itemize}
\item Compare your construction with a neighbour.
\end{enumerate}
\note{Do not identify conjunction elimination yet. Let participants describe projection in ordinary language.}
\end{frame}
\begin{frame}[plain]\titlepage\end{frame}
\begin{frame}{The route through the clinic}
\tableofcontents
\end{frame}
\subsection{Judgments}
\begin{frame}{A judgment records what is available}
\begin{definition}
The judgment $\seq{\Gamma}{A}$ says that $A$ has a derivation from assumptions in $\Gamma$.
\end{definition}
\begin{exampleblock}{Context}
$\Gamma=A\land B$ is a resource for proof search, not a truth-table row.
\end{exampleblock}
\end{frame}
\section{Build, then inspect}
\begin{frame}{Introduction tells us what to construct}
\begin{prooftree}
\AxiomC{$\Gamma\vdash A$}
\AxiomC{$\Gamma\vdash B$}
\RightLabel{$\land I$}
\BinaryInfC{$\Gamma\vdash A\land B$}
\end{prooftree}
\medskip
To finish with a conjunction, plan two subproofs.
\pause
\begin{alertblock}{Proof-search prompt}
Which component can you obtain first from the given assumption?
\end{alertblock}
\end{frame}
\begin{frame}{The completed proof has no hidden step}
\begin{prooftree}
\AxiomC{$A\land B$}
\RightLabel{$\land E_2$}\UnaryInfC{$B$}
\AxiomC{$A\land B$}
\RightLabel{$\land E_1$}\UnaryInfC{$A$}
\RightLabel{$\land I$}\BinaryInfC{$B\land A$}
\end{prooftree}
\smallskip
Read from leaves to root to check validity; read backwards to plan a proof.
\end{frame}
\begin{frame}{Turn the construction into a function}
\begin{theorem}[Conjunction exchange]
$A\land B\to B\land A$ is constructively derivable.
\end{theorem}
\begin{proof}
Assume $p:A\land B$. Project its second and first components, then pair them in that order. Discharge $p$.
\end{proof}
\[
 p\longmapsto (\pi_2p,\pi_1p).
\]
\end{frame}
\begin{frame}{Find the mistake in this tempting argument}
\begin{quote}
``Either $A$ or $B$ holds, so choose $A$.''
\end{quote}
\pause
\begin{alertblock}{Missing information}
A proof of $A\lor B$ does not generally supply a proof of $A$.
\end{alertblock}
\begin{itemize}
\item To eliminate a disjunction, handle \emph{both} cases.
\item Each branch must establish the same target.
\end{itemize}
\end{frame}
\section{Transfer}
\begin{frame}{Now change the target}
\begin{enumerate}[<+->]
\item Derive $A\land(B\land C)\to (A\land B)\land C$.
\item Mark where the same assumption is used more than once.
\item Explain why this exercise would change in a resource-sensitive logic.
\end{enumerate}
\bigskip
\textbf{Exit ticket:} describe introduction and elimination without using either word.
\end{frame}
\begin{frame}{Further reading}
\small
\begin{thebibliography}{9}
\bibitem{pfenning} Frank Pfenning. \emph{Natural Deduction}, 15-317: Constructive Logic, Carnegie Mellon University.
\href{https://www.cs.cmu.edu/~fp/courses/15317-f17/lectures/02-natded.pdf}{Course notes}.
\end{thebibliography}
The examples and proof-clinic sequence in this deck are original. The notation is standard intuitionistic natural deduction; no classical excluded-middle rule has been assumed.
\end{frame}
\end{document}
